\documentclass[aspectratio=169,xcolor=dvipsnames]{beamer}

\AtBeginSection[]
{
    \begin{frame}
        \frametitle{Visão Geral}
        \tableofcontents[currentsection]
    \end{frame}
}

\usefonttheme{professionalfonts} % using non standard fonts for beamer
\usefonttheme{serif} % default family is serif
\usepackage{fontspec}
\setmainfont{XCharter}[
    UprightFont    = *-Roman,
    BoldFont       = *-Bold,
    ItalicFont     = *-Italic,
    BoldItalicFont = *-BoldItalic,
]

\usepackage{hyperref}
\usepackage{graphicx} % Allows including images
\usepackage{booktabs} % Allows the use of \toprule, \midrule and \bottomrule in tables
\input{./slides/SimplePlusTheme/beamerthemeSimplePlus.sty}
\input{./slides/SimplePlusTheme/beamerinnerthemeSimplePlus.sty}
\input{./slides/SimplePlusTheme/beamerfontthemeSimplePlus.sty}
\input{./slides/SimplePlusTheme/beamercolorthemeSimplePlus.sty}

\usepackage{listings}
\setbeamercovered{transparent}

\usepackage{biblatex}
\addbibresource{slides/reference.bib}
\renewcommand*{\bibfont}{\footnotesize}
% Tamanho dos slides de "Leituras Recomendadas": maior que a bibliografia
% completa (\bibfont acima, que rende \tiny via a redefinição de \footnotesize)
% para destacar as leituras curadas.
\newcommand{\leiturasize}{\small}

\usepackage[brazilian ]{babel}
\usepackage{csquotes}
\usepackage{pgfplots}
\pgfplotsset{compat=1.18}

\usepackage{emoji}
\usepackage{amsmath}
\usepackage[xcharter,bigdelims,vvarbb]{newtxmath}
% newtxmath's operators font requests the fontspec family `XCharter(0)' in OT1;
% without these substitutions math digits and operator names silently fall back
% to Computer Modern (LaTeX Font Warning: Font shape `OT1/XCharter(0)/m/n' undefined).
\DeclareFontFamily{OT1}{XCharter(0)}{}
\DeclareFontShape{OT1}{XCharter(0)}{m}{n}{<->ssub * XCharter-TLF/m/n}{}
\DeclareFontShape{OT1}{XCharter(0)}{m}{it}{<->ssub * XCharter-TLF/m/it}{}
\DeclareFontShape{OT1}{XCharter(0)}{b}{n}{<->ssub * XCharter-TLF/b/n}{}
\DeclareFontShape{OT1}{XCharter(0)}{bx}{n}{<->ssub * XCharter-TLF/b/n}{}

\usepackage{bm}
\usepackage[table]{xcolor}
\usepackage{xcolor}
\usepackage{algpseudocode}

\usepackage{cancel}
\usepackage{ulem}

\usetikzlibrary{shapes}
\usetikzlibrary{arrows}
\usetikzlibrary {arrows.meta}
\usetikzlibrary{calc,positioning}
\usepgfplotslibrary{fillbetween}
\usetikzlibrary{angles,quotes}
\usetikzlibrary{tikzmark}


\definecolor{PastelRed}{HTML}{FFC1CC}
\definecolor{PastelBlue}{HTML}{C2F0FF}
\definecolor{PastelBlue2}{HTML}{3182BD}

\definecolor{PastelPink}{HTML}{FFCBE1}
\definecolor{PastelGreen}{HTML}{D6E5BD}
\definecolor{PastelYellow}{HTML}{F9E1A8}
\definecolor{PastelBlue}{HTML}{BCD8EC}
\definecolor{PastelPurple}{HTML}{DCCCEC}
\definecolor{PastelOrange}{HTML}{FFDAB4}


\renewcommand{\footnotesize}{\tiny}

\newcommand{\mespor}{
  \ifcase\month
    \or janeiro
    \or fevereiro
    \or março
    \or abril
    \or maio
    \or junho
    \or julho
    \or agosto
    \or setembro
    \or outubro
    \or novembro
    \or dezembro
  \fi
}
% \usepackage{csquotes}
% \usepackage{minted}
% \usemintedstyle{xcode}
\definecolor{vscLightBackground}{HTML}{FFFFFF}
\definecolor{vscLightText}{HTML}{000000}
\definecolor{vscLightGreen}{HTML}{008000}        % For Comments AND Numbers
\definecolor{vscLightRed}{HTML}{A31515}          % For Strings
\definecolor{vscLightTeal}{HTML}{267F99}
\definecolor{vscBlue}{HTML}{0000FF}% For main keywords (if, for, def...)
\definecolor{vscLightPytorch}{HTML}{800080}      % For PyTorch keywords (dark magenta)

\lstdefinestyle{vscode-light}{
    language=Python,
    backgroundcolor=\color{vscLightBackground},
    basicstyle=\ttfamily\scriptsize\color{vscLightText},
    keywordstyle=\color{vscLightTeal},
    commentstyle=\color{vscLightGreen},
    stringstyle=\color{vscLightRed},
    % Style for PyTorch keywords (using keyword group 2)
    keywordstyle=[2]\color{vscLightPytorch}, 
    morekeywords={self, True, False, None},
    morekeywords=[2]{
        torch, Tensor, device, dtype, to, nn, optim, autograd, utils, data,
        Module, Linear, Conv2d, ReLU, CrossEntropyLoss, Adam, SGD,
        Dataset, DataLoader, randn, zeros, ones, arange, cat, stack,
        view, reshape, sum, mean, max, min, matmul, no_grad, backward
    },
    keywordstyle=[3]\color{vscBlue}, 
    morekeywords=[3]{
        model, loss, y, y_hat, x , optimizer
    },
    frame=single,
    tabsize=1,
    showstringspaces=false,
    breaklines=true,
    captionpos=b,
    extendedchars=true,
    numbers=left,
}

\usepackage[cache=true]{minted}
% minted v3+ (TeX Live 2024+) auto-detects the output directory.
% minted v2 (Ubuntu apt TeX Live 2023) needs it set explicitly to
% match latexmkrc's $out_dir = 'build'.
\makeatletter
\@ifpackagelater{minted}{2024/01/01}{}{%
  \def\minted@outputdir{build/}%
}
\makeatother
\setminted{
    fontsize=\scriptsize,
    style=vs,
    breaklines=true,
    numbers=left,
    numbersep=5pt,
    frame=single,
}
\providecommand{\citeauthoryear}[1]{(\citeauthor{#1},~\citeyear{#1})}

\providecommand{\thetav}{\bm{\theta}}

% vetor coluna 2D com colchetes
\providecommand{\cvec}[2]{\begin{bmatrix} #1 \\ #2 \end{bmatrix}}

\providecommand{\varvector}[1]{\mathbf{#1}}
% instance
\providecommand{\X}{\varvector{X}}
\providecommand{\mlinstance}{\varvector{X}}
\providecommand{\mlinstancei}{\varvector{X}^{(i)}}
\providecommand{\mlinstancex}[1]{\varvector{X}^{(#1)}}

% target
\providecommand{\y}{\varvector{y}}
\providecommand{\mltarget}{\varvector{y}}
\providecommand{\mltargeti}{\varvector{y}^{(i)}}
\providecommand{\mltargetx}[1]{\varvector{y}^{(#1)}}

%hats
\providecommand{\yhat}{\hat{\varvector{y}}}
\providecommand{\mltargethat}{\hat{\varvector{y}}}
\providecommand{\mltargethati}{\hat{\varvector{y}}^{(i)}}
\providecommand{\mltargethatx}[1]{\hat{\varvector{y}}^{(#1)}}

% linear reg
\providecommand{\mllinearregression}{\mltargethat=\mlinstance\bm{\theta}}

\providecommand{\tikzarrowcolor}[3]{%
  \begin{tikzpicture}[remember picture, overlay]
    % Arrow body
    \path[fill=#3] ([shift={(#1,#2)}]current page.south west) 
      rectangle ++(0.3, -0.2);
    % Arrow head
    \path[fill=#3] ([shift={(#1+0.3,#2+0.1)}]current page.south west) -- 
         ++(0, -0.4) -- 
         ++(0.2, 0.2) -- 
         cycle;
  \end{tikzpicture}%
}

\providecommand{\tikzarrow}[2]{%
    \tikzarrowcolor{#1}{#2}{Red}
}



\DeclareMathOperator*{\argmax}{argmax}
\DeclareMathOperator*{\argmin}{argmin}



%----------------------------------------------------------------------------------------
%    TITLE PAGE
%----------------------------------------------------------------------------------------

\title{Deep Learning}
\subtitle{Inicialização de Pesos: Explosão e Dissipação de Gradientes}

\author{Lucas Silveira Kupssinskü}

\institute
{
    Machine Learning Theory and Applications Laboratory \\
    PUCRS
}
\date{\mespor de \the\year}

%----------------------------------------------------------------------------------------
%    PRESENTATION SLIDES
%----------------------------------------------------------------------------------------

\begin{document}

  \begin{frame}
    \titlepage
  \end{frame}

  \begin{frame}{Visão Geral}
    \tableofcontents
  \end{frame}

  %======================================================================
  \section{Por que inicialização importa}
  %======================================================================

  %--- Slide: recap das ativações ---
  \begin{frame}{Recap: nem toda ativação preserva o gradiente}
    \begin{columns}
      \begin{column}{.5\linewidth}
        \begin{itemize}
          \item Sigmoid e tanh \alert{saturam} para $|z|$ grande, com derivada
            tendendo a zero.
          \item ReLU e variantes preservam gradiente no lado positivo.
          \item A combinação \textbf{inicialização $\times$ ativação} controla a
            magnitude do sinal que se propaga em forward e em backward.
        \end{itemize}
      \end{column}
      \begin{column}{.5\linewidth}
        \centering
        \resizebox{\linewidth}{!}{
          \begin{tikzpicture}
            \begin{axis}[
              width=8cm,height=4.5cm,
              xlabel={$z$}, axis lines=middle,
              domain=-4:4, samples=120,
              ymin=-1.1, ymax=3.1,
              legend style={font=\tiny,at={(1.05,1)},anchor=north west}
            ]
              \addplot[thick,blue] {tanh(x)};
              \addlegendentry{tanh}
              \addplot[thick,red] {max(0,x)};
              \addlegendentry{ReLU}
            \end{axis}
          \end{tikzpicture}
        }
      \end{column}
    \end{columns}
  \end{frame}

  %--- Slide: papel da inicialização ---
  \begin{frame}{No início do treino, só os pesos definem o sinal}
    \begin{itemize}
      \item Antes da primeira atualização, a rede não foi ajustada aos dados.
      \item A magnitude das ativações em cada camada é determinada inteiramente pelos
        pesos iniciais e pelas funções de ativação.
      \item Se a magnitude das ativações crescer ou encolher exponencialmente
        camada após camada, o gradiente faz o mesmo.
    \end{itemize}
  \end{frame}

  %--- Slide: vanishing/exploding ---
  \begin{frame}{Vanishing e exploding gradients}
    \begin{itemize}
      \item \textbf{Vanishing gradient}: gradientes encolhem a cada camada, camadas
        iniciais não recebem sinal de aprendizagem.
      \item \textbf{Exploding gradient}: gradientes crescem a cada camada, atualizações
        ficam instáveis numericamente.
      \item Ambos os casos quebram a descida de gradiente.
      \item A inicialização \alert{certa} mantém a variância das pré-ativações
        aproximadamente constante ao longo das camadas.
    \end{itemize}
  \end{frame}

  %======================================================================
  \section{Estudo de caso: MNIST com tanh}
  %======================================================================

  %--- Slide: setup ---
  \begin{frame}{Setup do estudo de caso}
    \begin{columns}
      \begin{column}{.4\linewidth}
        \centering
        \includegraphics[width=\linewidth,height=.55\textheight,keepaspectratio]{slides/img/01-mlp/init-mnist.png}
      \end{column}
      \begin{column}{.6\linewidth}
        \begin{itemize}
          \item Dataset: MNIST, dígitos manuscritos.
          \item Arquitetura: MLP com 5 camadas, todas com 300 unidades tanh,
            saída softmax com 10 classes.
          \item Tamanho do \textit{batch}: 100 imagens.
          \item Vamos olhar o histograma das ativações $A^{[1]},\ldots,A^{[4]}$
            em cada camada oculta.
        \end{itemize}
      \end{column}
    \end{columns}
    \vfill
    \centering
    \resizebox{.95\linewidth}{!}{%
    \begin{tikzpicture}[>=stealth, every node/.style={font=\small}]
      \node (x) at (0, 0) {$X = A^{[0]}$};
      \node[below=0.15cm of x, font=\scriptsize, align=center]
        {batch de 100 imagens $28\times28$\\shape $(100, 784)$};
      \node[circle, draw, minimum size=1.1cm] (h1) at (2.6, 0) {tanh};
      \node[circle, draw, minimum size=1.1cm] (h2) at (5.2, 0) {tanh};
      \node[circle, draw, minimum size=1.1cm] (h3) at (7.8, 0) {tanh};
      \node[circle, draw, minimum size=1.1cm] (h4) at (10.4, 0) {tanh};
      \node[circle, draw, minimum size=1.1cm, font=\scriptsize] (out) at (13.2, 0) {softmax};
      \node (y) at (15.8, 0) {$\hat{y} = A^{[5]}$};
      \node[below=0.15cm of y, font=\scriptsize, align=center]
        {probabilidades das 10 classes\\shape $(100, 10)$};
      \draw[->] (x)  -- (h1);
      \draw[->] (h1) -- node[above]{$A^{[1]}$} node[below, font=\scriptsize]{$(100, 300)$} (h2);
      \draw[->] (h2) -- node[above]{$A^{[2]}$} node[below, font=\scriptsize]{$(100, 300)$} (h3);
      \draw[->] (h3) -- node[above]{$A^{[3]}$} node[below, font=\scriptsize]{$(100, 300)$} (h4);
      \draw[->] (h4) -- node[above]{$A^{[4]}$} node[below, font=\scriptsize]{$(100, 300)$} (out);
      \draw[->] (out) -- (y);
      \node[below=0.15cm of h1, font=\scriptsize, align=center]
        {300 unidades\\$\theta^{[1]}\!: (300, 784)$};
      \node[below=0.15cm of h2, font=\scriptsize, align=center]
        {300 unidades\\$\theta^{[2]}\!: (300, 300)$};
      \node[below=0.15cm of h3, font=\scriptsize, align=center]
        {300 unidades\\$\theta^{[3]}\!: (300, 300)$};
      \node[below=0.15cm of h4, font=\scriptsize, align=center]
        {300 unidades\\$\theta^{[4]}\!: (300, 300)$};
      \node[below=0.15cm of out, font=\scriptsize, align=center]
        {10 unidades\\$\theta^{[5]}\!: (10, 300)$};
    \end{tikzpicture}%
    }
  \end{frame}

  %--- Slide: theta = 0 ---
  \begin{frame}{Tentativa 1: $\theta = 0$}
    \begin{columns}
      \begin{column}{.4\linewidth}
        \begin{itemize}
          \item Pesos iniciais zerados.
          \item Toda unidade computa a mesma pré-ativação.
          \item Todas as ativações colapsam em 0.
          \item Por simetria, qualquer atualização afeta todas as unidades igualmente:
            a rede nunca quebra a simetria.
        \end{itemize}
      \end{column}
      \begin{column}{.6\linewidth}
        \centering
        \includegraphics[width=\linewidth]{slides/img/01-mlp/init-hist-zero.png}
        \footnotesize Histogramas de $A^{[1]},\ldots,A^{[4]}$ com $\theta=0$.
      \end{column}
    \end{columns}
  \end{frame}

  %--- Slide: theta ~ N(0,1) ---
  \begin{frame}{Tentativa 2: $\theta \sim \mathcal{N}(0,1)$}
    \begin{columns}
      \begin{column}{.4\linewidth}
        \begin{itemize}
          \item Pesos com variância 1.
          \item As pré-ativações de cada camada explodem em magnitude.
          \item A tanh \alert{satura} em $\pm 1$.
          \item Em camadas internas, quase todas as unidades estão na região
            de derivada nula.
        \end{itemize}
      \end{column}
      \begin{column}{.6\linewidth}
        \centering
        \includegraphics[width=\linewidth]{slides/img/01-mlp/init-hist-normal.png}
        \footnotesize Massa concentrada em $\pm 1$: zona morta da tanh.
      \end{column}
    \end{columns}
  \end{frame}

  %--- Slide: theta ~ U(-1, 1) ---
  \begin{frame}{Tentativa 3: $\theta \sim \mathcal{U}(-1,1)$}
    \begin{columns}
      \begin{column}{.4\linewidth}
        \begin{itemize}
          \item Variância $\sigma_\theta^2 = 1/3$, menor que o caso anterior.
          \item Mesmo assim, as ativações \alert{encolhem} a cada camada.
          \item Em $A^{[4]}$ a distribuição é uma agulha em zero.
          \item Sintoma de \textit{vanishing} em forward.
        \end{itemize}
      \end{column}
      \begin{column}{.6\linewidth}
        \centering
        \includegraphics[width=\linewidth]{slides/img/01-mlp/init-hist-uniform.png}
        \footnotesize Distribuição encolhe a cada camada.
      \end{column}
    \end{columns}
  \end{frame}

  %--- Slide: conclusão da seção ---
  \begin{frame}{Precisamos calibrar $\sigma_\theta^2$ por camada}
    \begin{itemize}
      \item Inicializações ad-hoc falham de formas opostas dependendo da escala.
      \item A escala certa depende do número de unidades de entrada de cada camada
        e da função de ativação utilizada.
      \item Visualização interativa que motiva a derivação:
        \begin{center}
          \href{https://lucaskup.github.io/deep-learning-course/weight-init/}{%
            \textit{\color{blue}\underline{Demo interativo: Inicialização de Pesos}}
          }
        \end{center}
    \end{itemize}
  \end{frame}

  %======================================================================
  \section{Forward pass: média da pré-ativação}
  %======================================================================

  %--- Slide: definição ---
  \begin{frame}{Definição da pré-ativação por camada}
    Cada exemplo é um vetor linha e o peso multiplica a entrada pela direita:
    \[
      z^{(l)} = \varphi\!\left(z^{(l-1)}\right)\theta^{(l)\top} + b^{(l)}
      \qquad\Longrightarrow\qquad
      z^{(l)}_i = b^{(l)}_i + \sum_{j=1}^{n^{(l-1)}} \varphi\!\left(z^{(l-1)}_j\right)\theta^{(l)}_{ij}
    \]
    \begin{itemize}
      \item $z^{(l)}$: pré-ativação da camada $l$, vetor linha com $n^{(l)}$ componentes.
      \item $b^{(l)}$: bias, \alert{inicializado em 0}.
      \item $\theta^{(l)}$: pesos, matriz de shape $(n^{(l)}, n^{(l-1)})$,
        com $\theta^{(l)}_{ij} \sim \mathcal{N}(0,\sigma_\theta^2)$.
      \item $\varphi$: ativação (tomamos ReLU como referência).
    \end{itemize}
    \vfill
    \begin{block}{Hipóteses}
      $\theta^{(l)}_{ij}$ independentes entre si e independentes de
      $\varphi(z^{(l-1)}_j)$. A distribuição de $z^{(l-1)}_j$ é simétrica em zero.
    \end{block}
  \end{frame}

  %--- Slide: linearidade da esperança ---
  \begin{frame}{Esperança: separar bias do somatório}
    Aplicando a linearidade da esperança:
    \[
      E[z^{(l)}_i]
        = E[b^{(l)}_i] + E\!\left[\sum_{j=1}^{n^{(l-1)}} \varphi(z^{(l-1)}_j)\,\theta^{(l)}_{ij}\right].
    \]
    \begin{itemize}
      \item Como inicializamos $b^{(l)}_i = 0$, $E[b^{(l)}_i] = 0$.
      \item Sobra a esperança do somatório, que reescrevemos como soma das esperanças.
    \end{itemize}
    \[
      E[z^{(l)}_i] = \sum_{j=1}^{n^{(l-1)}} E\!\left[\varphi(z^{(l-1)}_j)\,\theta^{(l)}_{ij}\right].
    \]
  \end{frame}

  %--- Slide: independência ---
  \begin{frame}{Esperança: independência $\varphi\perp\theta$}
    Como $\varphi(z^{(l-1)}_j)$ é independente de $\theta^{(l)}_{ij}$,
    \[
      E\!\left[\varphi(z^{(l-1)}_j)\,\theta^{(l)}_{ij}\right]
        = E[\varphi(z^{(l-1)}_j)]\,E[\theta^{(l)}_{ij}].
    \]
    \begin{itemize}
      \item Por construção, $\theta^{(l)}_{ij}\sim\mathcal{N}(0,\sigma_\theta^2)$, logo
        $E[\theta^{(l)}_{ij}] = 0$.
      \item O produto inteiro é zero, independentemente de $\varphi$.
    \end{itemize}
    \[
      E[z^{(l)}_i] = \sum_{j=1}^{n^{(l-1)}} E[\varphi(z^{(l-1)}_j)] \cdot 0 = 0.
    \]
  \end{frame}

  %--- Slide: resultado parcial ---
  \begin{frame}{Resultado intermediário}
    \begin{block}{Média das pré-ativações é zero}
      \[
        E[z^{(l)}_i] = 0 \quad\text{para toda camada } l.
      \]
    \end{block}
    \vfill
    \begin{itemize}
      \item Vamos usar este resultado para calcular a variância em seguida.
      \item Como $E[z^{(l)}_i]=0$, temos $\mathrm{Var}[z^{(l)}_i] = E[(z^{(l)}_i)^2]$.
    \end{itemize}
  \end{frame}

  %======================================================================
  \section{Forward pass: variância da pré-ativação}
  %======================================================================

  %--- Slide: definição da variância ---
  \begin{frame}{Variância: ponto de partida}
    \[
      \sigma^2_{z^{(l)}}
        = E\!\left[(z^{(l)}_i)^2\right] - \cancelto{0}{\big(E[z^{(l)}_i]\big)^2}
        = E\!\left[\Big(b^{(l)}_i + \sum_{j} \varphi(z^{(l-1)}_j)\,\theta^{(l)}_{ij}\Big)^2\right].
    \]
    \begin{itemize}
      \item Como $b^{(l)}_i = 0$, expandimos sem o termo do bias.
      \item Cross-terms entre pesos diferentes desaparecem por independência
        e $E[\theta]=0$.
    \end{itemize}
    \[
      \sigma^2_{z^{(l)}}
        = \sum_{j=1}^{n^{(l-1)}} E\!\left[\varphi\!\big(z^{(l-1)}_j\big)^2\right]\,
                                 E\!\left[\big(\theta^{(l)}_{ij}\big)^2\right].
    \]
  \end{frame}

  %--- Slide: identifica sigma_theta^2 ---
  \begin{frame}{Identificando $\sigma_\theta^2$}
    Como $\theta^{(l)}_{ij}\sim\mathcal{N}(0,\sigma_\theta^2)$ e $E[\theta]=0$:
    \[
      E\!\left[\big(\theta^{(l)}_{ij}\big)^2\right]
        = \mathrm{Var}[\theta^{(l)}_{ij}] = \sigma_\theta^2.
    \]
    Substituindo:
    \[
      \sigma^2_{z^{(l)}}
        = \sum_{j=1}^{n^{(l-1)}} E\!\left[\varphi\!\big(z^{(l-1)}_j\big)^2\right]\sigma_\theta^2
        = \sigma_\theta^2 \sum_{j=1}^{n^{(l-1)}} E\!\left[\varphi\!\big(z^{(l-1)}_j\big)^2\right].
    \]
    \begin{itemize}
      \item Falta calcular $E[\varphi(z)^2]$ para $\varphi$ = ReLU.
    \end{itemize}
  \end{frame}

  %--- Slide: ReLU sobre simétrica ---
  \begin{frame}{Para ReLU sobre $z$ simétrica em zero}
    Vamos provar (próximo slide) que, se $z$ é simétrica em zero com variância $\sigma_z^2$,
    \[
      E[\mathrm{ReLU}(z)^2] = \frac{\sigma_z^2}{2}.
    \]
    \begin{itemize}
      \item ReLU corta os valores negativos, mantendo metade da massa.
      \item A simetria garante que a metade positiva tem o mesmo segundo momento da
        distribuição inteira, dividido por 2.
    \end{itemize}
  \end{frame}

  %--- Slide: Junta tudo ---
  \begin{frame}{Variância da pré-ativação em função de $\sigma_\theta^2$ e $\sigma^2_{z^{(l-1)}}$}
    Aplicando o resultado anterior, $E[\varphi(z^{(l-1)}_j)^2] = \sigma^2_{z^{(l-1)}}/2$
    para toda unidade $j$:
    \[
      \sigma^2_{z^{(l)}}
        = \sigma_\theta^2 \sum_{j=1}^{n^{(l-1)}} \frac{\sigma^2_{z^{(l-1)}}}{2}
        = \frac{n^{(l-1)}}{2}\,\sigma_\theta^2\,\sigma^2_{z^{(l-1)}}.
    \]
    \vfill
    \begin{block}{Recorrência para a variância}
      \[
        \boxed{\;\sigma^2_{z^{(l)}} = \tfrac{n^{(l-1)}}{2}\,\sigma_\theta^2\,\sigma^2_{z^{(l-1)}}\;}
      \]
    \end{block}
    \begin{itemize}
      \item Se a constante $\tfrac{n^{(l-1)}}{2}\,\sigma_\theta^2$ for maior que 1, exploding.
      \item Se for menor que 1, vanishing. Queremos exatamente 1.
    \end{itemize}
  \end{frame}

  %======================================================================
  \section{Exercício teórico}
  %======================================================================

  %--- Slide: enunciado ---
  \begin{frame}{Exercício: segundo momento de ReLU}
    \begin{exampleblock}{Enunciado}
      Seja $a$ uma variável aleatória contínua, simétrica em torno de $E[a]=0$, com
      variância $\mathrm{Var}[a] = \sigma^2$. Mostre que
      \[
        E\!\left[\mathrm{ReLU}(a)^2\right] = \frac{\sigma^2}{2}.
      \]
    \end{exampleblock}
    \vfill
    \begin{itemize}
      \item Lembre que $\mathrm{ReLU}(a) = \max(0,a)$.
      \item Use $E[a^2] = \mathrm{Var}[a] + (E[a])^2 = \sigma^2$.
      \item Use a simetria de $a$ em torno de zero.
    \end{itemize}
  \end{frame}

  %--- Slide: solução ---
  \begin{frame}{Solução}
    Pela definição de $\mathrm{ReLU}$:
    \[
      E[\mathrm{ReLU}(a)^2] = \int_{-\infty}^{\infty} \max(0,a)^2 p(a)\,da
        = \int_{0}^{\infty} a^2 p(a)\,da.
    \]
    Pela simetria de $p$ em torno de zero, $\int_{-\infty}^{0} a^2 p(a)\,da = \int_{0}^{\infty} a^2 p(a)\,da$.
    Logo,
    \[
      E[a^2] = 2 \int_{0}^{\infty} a^2 p(a)\,da
        \quad\Longrightarrow\quad
      \int_{0}^{\infty} a^2 p(a)\,da = \frac{E[a^2]}{2} = \frac{\sigma^2}{2}.
    \]
    \[
      \boxed{E[\mathrm{ReLU}(a)^2] = \frac{\sigma^2}{2}.}
    \]
  \end{frame}

  %======================================================================
  \section{Inicialização He / Kaiming}
  %======================================================================

  %--- Slide: condição de invariância ---
  \begin{frame}{Condição: variância igual a cada camada}
    Queremos $\sigma^2_{z^{(l)}} = \sigma^2_{z^{(l-1)}}$ para todo $l$. Da recorrência,
    \[
      \sigma^2_{z^{(l-1)}} = \tfrac{n^{(l-1)}}{2}\,\sigma_\theta^2\,\sigma^2_{z^{(l-1)}}
      \;\Longrightarrow\;
      \tfrac{n^{(l-1)}}{2}\,\sigma_\theta^2 = 1.
    \]
    \vfill
    \begin{block}{He / Kaiming}
      \[
        \boxed{\;\sigma_\theta^2 = \dfrac{2}{n^{(l-1)}}\;}
        \qquad
        \theta^{(l)}_{ij} \sim \mathcal{N}\!\left(0,\;\tfrac{2}{n^{(l-1)}}\right).
      \]
    \end{block}
    \begin{itemize}
      \item $n^{(l-1)}$ é o número de \textit{fan-in} da camada (entradas).
      \item Padrão de fato para ReLU e GELU.
    \end{itemize}
  \end{frame}

  %--- Slide: He paper ---
  \begin{frame}{Origem da inicialização He\footfullcite{he2015delving}}
    \begin{itemize}
      \item Apresentada em 2015 por Kaiming He et al. para redes muito profundas com ReLU.
      \item Permitiu treinar redes com mais de 30 camadas convolucionais sem
        \textit{batch norm}.
      \item Variante uniforme: $\theta\sim\mathcal{U}\!\left(-\sqrt{6/n},\;\sqrt{6/n}\right)$,
        que tem a mesma variância $2/n$.
    \end{itemize}
  \end{frame}

  %--- Slide: Xavier / Glorot ---
  \begin{frame}{Inicialização Xavier (Glorot)\footfullcite{glorot2010xavier}}
    \begin{itemize}
      \item Antes de He, Glorot e Bengio (2010) fizeram o mesmo argumento para
        \textbf{tanh} ao redor de zero.
      \item Para tanh, $\varphi(z)\approx z$ perto da origem, então $E[\varphi(z)^2]\approx \sigma_z^2$
        sem o fator $1/2$.
      \item Isso leva a $\sigma_\theta^2 = 1/n^{(l-1)}$.
      \[
        \boxed{\;\sigma_\theta^2 = \dfrac{1}{n^{(l-1)}}\;}
        \qquad
        \theta^{(l)}_{ij} \sim \mathcal{N}\!\left(0,\;\tfrac{1}{n^{(l-1)}}\right).
      \]
      \item Padrão para sigmoid, tanh e softmax.
    \end{itemize}
  \end{frame}

  %======================================================================
  \section{Voltando ao MNIST}
  %======================================================================

  %--- Slide: He resolve o problema ---
  \begin{frame}{Mesmo MLP, agora com inicialização He}
    \begin{columns}
      \begin{column}{.6\linewidth}
        \centering
        \includegraphics[width=\linewidth,height=0.7\textheight,keepaspectratio]{slides/img/01-mlp/init-hist-he.png}
      \end{column}
      \begin{column}{.4\linewidth}
        \begin{itemize}
          \item Histogramas das ativações $A^{[1]},\ldots,A^{[4]}$ ficam estáveis
            ao longo das camadas.
          \item A magnitude não cresce nem encolhe, e a tanh não satura.
          \item A rede consegue treinar.
        \end{itemize}
      \end{column}
    \end{columns}
  \end{frame}

  %--- Slide: comparação visual ---
  \begin{frame}{Comparação visual das quatro inicializações}
    \begin{columns}
      \begin{column}{.5\linewidth}
        \centering
        \includegraphics[width=.9\linewidth,height=.36\textheight,keepaspectratio]{slides/img/01-mlp/init-hist-zero.png}\\
        \footnotesize $\theta = 0$\\[0.3em]
        \includegraphics[width=.9\linewidth,height=.36\textheight,keepaspectratio]{slides/img/01-mlp/init-hist-uniform.png}\\
        \footnotesize $\theta\sim\mathcal{U}(-1,1)$
      \end{column}
      \begin{column}{.5\linewidth}
        \centering
        \includegraphics[width=.9\linewidth,height=.36\textheight,keepaspectratio]{slides/img/01-mlp/init-hist-normal.png}\\
        \footnotesize $\theta\sim\mathcal{N}(0,1)$\\[0.3em]
        \includegraphics[width=.9\linewidth,height=.36\textheight,keepaspectratio]{slides/img/01-mlp/init-hist-he.png}\\
        \footnotesize He: $\theta\sim\mathcal{N}(0,2/n)$
      \end{column}
    \end{columns}
  \end{frame}

  %======================================================================
  \section{Forward + Backward}
  %======================================================================

  %--- Slide: backward simétrico ---
  \begin{frame}{O argumento simétrico para o gradiente}
    \begin{itemize}
      \item O mesmo raciocínio aplicado à variância dos gradientes em \textit{backward}
        leva a uma condição análoga.
      \item Em vez de $n^{(l-1)}$ (\textit{fan-in}), aparece $n^{(l)}$ (\textit{fan-out}):
        $\sigma_\theta^2 = 2/n^{(l)}$.
      \item Os dois lados nem sempre são compatíveis (camadas com formas diferentes).
    \end{itemize}
  \end{frame}

  %--- Slide: fan_avg ---
  \begin{frame}{Compromisso fan\_avg}
    Glorot propôs balancear forward e backward usando a média de
    $n^{(l-1)}$ e $n^{(l)}$:
    \[
      \boxed{\;\sigma_\theta^2 = \dfrac{2}{n^{(l-1)} + n^{(l)}}\;\text{ (Xavier)}\;}
      \qquad
      \boxed{\;\sigma_\theta^2 = \dfrac{4}{n^{(l-1)} + n^{(l)}}\;\text{ (He, fan\_avg)}\;}
    \]
    \begin{itemize}
      \item Conhecido como \textit{fan\_avg} em frameworks.
      \item Em camadas onde $n^{(l-1)} = n^{(l)}$, recupera as fórmulas originais.
    \end{itemize}
    \vspace{0.3em}
    {\footnotesize\color{gray!50!black}A versão He original (He et al.\ 2015) usa apenas \textit{fan\_in} ($2/n^{(l-1)}$). A média \textit{fan\_avg} vem de Glorot/Xavier e é o default do \texttt{Dense} no Keras (\texttt{glorot\_uniform}); o PyTorch usa \textit{fan\_in} por padrão (ver próximo slide).}
  \end{frame}

  %======================================================================
  \section{Diretrizes práticas}
  %======================================================================

  %--- Slide: mapa ativação -> init ---
  \begin{frame}{Qual inicialização usar?}
    \centering
    \begin{tabular}{@{}lll@{}}
      \toprule
      Ativação principal & Inicialização recomendada & Variância $\sigma_\theta^2$ \\
      \midrule
      ReLU, Leaky ReLU, ELU & He / Kaiming           & $2/n^{(l-1)}$            \\
      GELU                 & He / Kaiming           & $2/n^{(l-1)}$            \\
      Sigmoid, tanh        & Xavier / Glorot        & $1/n^{(l-1)}$            \\
      Softmax (saída)      & Xavier / Glorot        & $1/n^{(l-1)}$            \\
      \bottomrule
    \end{tabular}
    \vfill
    \begin{itemize}
      \item Bias sempre inicializado em 0 (ou em valor pequeno positivo para ReLU,
        evita unidades mortas no início).
    \end{itemize}
  \end{frame}

  %--- Slide: PyTorch ---
  \begin{frame}{Em PyTorch}
    \begin{itemize}
      \item He (Kaiming): \texttt{torch.nn.init.kaiming\_normal\_(tensor, mode='fan\_in')}
      \item Xavier (Glorot): \texttt{torch.nn.init.xavier\_normal\_(tensor)}
      \item Padrão dos módulos \texttt{nn.Linear} e \texttt{nn.Conv2d}: \alert{Kaiming uniform}
        com \texttt{fan\_in}, $a=\sqrt{5}$. Bom o suficiente na maioria dos casos.
      \item Para arquiteturas grandes (transformers, ResNet), inicializações específicas
        por bloco (DeepNorm, MUP) podem dar resultados ainda melhores.
    \end{itemize}
  \end{frame}

  %--- Slide: leituras ---
  \section*{Referências e Leituras}
  \begin{frame}{Leituras Recomendadas}
    \leiturasize
    \sloppy
    \begin{itemize}
      \item Capítulo 7 de \fullcite{prince2023understanding}
      \item Xavier (Glorot e Bengio, 2010): \citeauthoryear{glorot2010xavier}
      \item Kaiming (He et al., 2015): \citeauthoryear{he2015delving}
      \item Visualização interativa:
        \href{https://lucaskup.github.io/deep-learning-course/weight-init/}{%
          \textit{lucaskup.github.io/deep-learning-course/weight-init}}
    \end{itemize}
  \end{frame}

  \begin{frame}[allowframebreaks]{Referências}
    \sloppy
    \printbibliography[heading=none]
  \end{frame}

\end{document}
